% AnoGen / GenFSDiff briefing deck.
%
%   make            rebuild figures and compile
%   make slides     compile only (figures unchanged)
%
% Figures come from slides/141026/make_figs.py, which reads the run outputs in
% results/. Every number below is sourced in docs/ESA_RUN_2026-10-06.md §11 and
% docs/MODEL_CONTRAST_HEAD.md; re-measure the morphology rows with
% scripts/measure_morphology.py.

% 9pt: a 16x9cm Beamer page is small relative to the type, and these frames are
% dense. Raise it if you cut content.
\documentclass[aspectratio=169,9pt,xcolor=table]{beamer}

\usepackage{amsmath,amssymb,booktabs,graphicx}
\usepackage[most]{tcolorbox}
\usepackage{tikz}

\usetheme{default}
\setbeamertemplate{navigation symbols}{}

% ---------------------------------------------------------------- palette
\definecolor{ink}{HTML}{16304A}      % body text, headings
\definecolor{accent}{HTML}{1F6FB4}   % our emphasis / "generated"
\definecolor{good}{HTML}{1A7A4C}     % results in our favour
\definecolor{bad}{HTML}{9B1C1C}      % caveats / "real"
\definecolor{mute}{HTML}{5C6B7A}     % secondary text
\definecolor{hair}{HTML}{D6DDE4}     % rules
\definecolor{band}{HTML}{F2F6FA}     % table striping

\newcommand{\gd}[1]{\textcolor{good}{\textbf{#1}}}
\newcommand{\bd}[1]{\textcolor{bad}{#1}}
\newcommand{\mt}[1]{\textcolor{mute}{#1}}

% ------------------------------------------------------------------ theme
\setbeamercolor{normal text}{fg=ink}
\setbeamercolor{frametitle}{fg=ink}
\setbeamercolor{itemize item}{fg=accent}
\setbeamercolor{itemize subitem}{fg=accent}
\setbeamerfont{frametitle}{size=\Large,series=\bfseries}
\setbeamerfont{framesubtitle}{size=\small}
\setbeamercolor{framesubtitle}{fg=mute}
\setbeamertemplate{itemize item}{\raisebox{0.15ex}{\scriptsize$\bullet$}}
\setbeamertemplate{itemize subitem}{\raisebox{0.15ex}{\tiny$\bullet$}}
\setlength{\leftmargini}{1.1em}
\setlength{\leftmarginii}{1.1em}

% Left-aligned frame title with a subtitle and a hairline under it.
% Beamer already indents this template to the body margin, so \hsize here is the
% body width (14cm) and no extra \hspace is wanted.
\makeatletter
\setbeamertemplate{frametitle}{%
  \vskip0.3em%
  \begin{minipage}{\hsize}%
    \raggedright
    {\usebeamerfont{frametitle}\usebeamercolor[fg]{frametitle}\insertframetitle}%
    \ifx\insertframesubtitle\@empty\else
      \\[0.1em]{\usebeamerfont{framesubtitle}%
        \usebeamercolor[fg]{framesubtitle}\insertframesubtitle}%
    \fi
  \end{minipage}%
  \vskip0.3em%
  {\color{hair}\rule{\hsize}{0.4pt}}%
  \vskip-0.3em%
}
\makeatother

% A box for an equation or a definition. #1 = title, #2 = colour.
\newtcolorbox{eqbox}[2][accent]{%
  colback=#1!5, colframe=#1!35, boxrule=0.5pt, arc=2pt,
  left=5pt, right=5pt, top=2pt, bottom=2pt,
  title={\color{ink}\bfseries\small #2}, coltitle=ink,
  fonttitle=\bfseries, colbacktitle=#1!5, titlerule=0pt,
  attach title to upper=\par\vspace{2pt}}

\newcommand{\tablestripe}{\rowcolors{2}{band}{white}}

% ------------------------------------------------------------------ meta
\title{AnoGen / GenFSDiff}
\subtitle{Few-shot steered diffusion for satellite-telemetry anomalies}
\date{14 October 2026}

\begin{document}

% ======================================================= 1. title
\begin{frame}[plain]
  \vspace{0.3cm}
  {\usebeamerfont{title}\Huge\bfseries\color{ink}GenFSDiff}\\[0.45em]
  {\large\color{accent}Few-shot steered diffusion for satellite-telemetry anomalies}\\[0.35em]
  {\small\color{mute}ESA-ADB Mission 1, channels 41--46}

  \vspace{0.45cm}
  {\color{hair}\rule{\textwidth}{0.4pt}}
  \vspace{0.25cm}

  \begin{itemize}\setlength\itemsep{0.35em}
    \item \textbf{Problem.} Seven years of telemetry per channel hold only \textbf{22}
          labelled anomaly events. \\ Too few anomalies to train a reasonable detector.
    \item \textbf{Approach.} A generative model to generate an anomaly of a chosen kind (+size and position).
    \item \textcolor{good}{\textbf{Result.} A detector trained on real anomalies plus generated ones improves its performance.}
  \end{itemize}

  \vfill
  {\footnotesize\color{mute}14 October 2026}
\end{frame}

% ======================================================= 2. the data
\begin{frame}{The data: ESA-ADB Mission 1}{}

  \vspace{0.3cm}
  \centering
  \tablestripe
  \small\renewcommand{\arraystretch}{1.2}
  \begin{tabular}{@{}l l r r@{}}
    \toprule
    the four kinds we generate & what it looks like & windows & events \\
    \midrule
    Point / Global            & one sample far outside the normal range      &  18 &  3 \\
    level shift               & the signal steps to a new level and stays    &   6 &  1 \\
    global subsequence        & a long stretch outside the normal range      & 636 & 17 \\
    local subsequence         & a stretch of normal values, wrong for the context & 12 & 2 \\
    \bottomrule
  \end{tabular}

  \vspace{0.18cm}
  \includegraphics[width=\textwidth]{figs/kinds.pdf}

  \vspace{0.1cm}
  \raggedright
  {\scriptsize\color{mute}}
\end{frame}

% ======================================================= 3. diffusion primer
\begin{frame}{What are diffusion models?}{The framework behind e.g. image generators;
  we apply it to telemetry}

  \centering
  \begin{tikzpicture}
    \node[inner sep=0] (strip) {\includegraphics[width=9.8cm]{figs/diffusion_strip.pdf}};
    % reverse process, above the strip
    \draw[->,thick,accent] ([yshift=5mm]strip.north east) -- ([yshift=5mm]strip.north west);
    \node[above,accent,font=\small] at ([yshift=6mm]strip.north)
      {reverse: a trained neural network $\varepsilon_\theta$ removes it one step at a time
      };
    % forward process, below
    \draw[->,thick,mute!60] ([yshift=-7mm]strip.south west) -- ([yshift=-7mm]strip.south east);
    \node[below,mute,font=\small] at ([yshift=-7.5mm]strip.south)
      {forward: add Gaussian noise --- fixed, nothing is learned};
    \node[below left,font=\scriptsize] at ([xshift=10mm]strip.south west) {$x_0$\ \ data};
    \node[below right,font=\scriptsize] at ([xshift=-10mm]strip.south east)
      {$x_T$\ \ pure noise};
  \end{tikzpicture}

  \vspace{0.35cm}

  \begin{columns}[T]
    \begin{column}{0.52\textwidth}
      \begin{eqbox}{Forward process $q(x_t \mid x_0)$,\\ model's guess of the original $\hat x_0$:}
        \small
        \vspace{-0.6em}
        \begin{align*}
          q(x_t \mid x_0) &= \mathcal{N}\bigl(\sqrt{\bar\alpha_t}\,x_0,\;
                             (1-\bar\alpha_t)I\bigr)\\[0.1em]
          \hat x_0 &= \bigl(x_t - \sqrt{1-\bar\alpha_t}\,
                      \varepsilon_\theta(x_t,t)\bigr)\big/\sqrt{\bar\alpha_t}
        \end{align*}
        \vspace{-1.1em}
      \end{eqbox}
    \end{column}
    \begin{column}{0.45\textwidth}
      \begin{itemize}\setlength\itemsep{0.25em}\small
        \item $\bar\alpha_t$: a fixed schedule \\--- how much of the
              original signal is left\\ at step $t$.
        \item $\varepsilon_\theta$: trained to predict the noise that was added.
      \end{itemize}
    \end{column}
  \end{columns}
\end{frame}

% ======================================================= 4. sampling
\begin{frame}{How a (nominal) sample is generated}{Start from noise, predict the added noise, and remove it. Repeat.}

  \begin{eqbox}[ink]{The loop, as implemented}
    \renewcommand{\arraystretch}{1.6}
    \centering
    \resizebox{0.85\linewidth}{!}{%
    \begin{tabular}{@{}l@{\hspace{2em}}l@{}}
      $x_T \sim \mathcal{N}(0,\,I)$
        & \mt{draw pure noise} \\
      \multicolumn{2}{@{}l@{}}{\textbf{for} $t = T\!-\!1,\;\dots,\;0$:} \\
      \quad $\varepsilon \leftarrow \varepsilon_\theta(x_t,\,t,\,c)$
        & \mt{one network call: the noise it sees in $x_t$} \\
      \quad $\hat x_0 \leftarrow \bigl(x_t - \sqrt{1-\bar\alpha_t}\,\varepsilon\bigr)
             \big/\sqrt{\bar\alpha_t}$
        & \mt{undo the forward process: guess of $x_0$} \\
      \quad $x_{t-1} \leftarrow \sqrt{\bar\alpha_{t-1}}\,\hat x_0
             + \sqrt{1-\bar\alpha_{t-1}}\,\varepsilon$
        & \mt{\emph{same} noise back, at level $t-1$} \\
      \quad $t \leftarrow t-1$
        & \mt{go to the next step} \\
      \textbf{return} $x_0$
    \end{tabular}}
    \vspace{-0.1em}
  \end{eqbox}
\end{frame}

% ======================================================= 5. our setting
\begin{frame}{Our setting: Steer the sample to generate an anomaly}{}

  \begin{eqbox}{One sampling step with the steering edit $\mathbf{E}$}
    \vspace{-0.3em}
    \[
      \color{ink}\boldsymbol{x_{t-1} \;=\; \sqrt{\bar\alpha_{t-1}}\,
      \bigl(\hat x_0 - \mathbf{E}\bigr) \;+\; \sqrt{1-\bar\alpha_{t-1}}\,\varepsilon}
    \]
    \vspace{-0.9em}
  \end{eqbox}
  \begin{center}
  $\mathbf{E}$ is the only thing we add --- the denoiser itself is never retrained.
  \end{center}
  \vspace{0.1cm}
  {\color{accent}

    \vspace{0.1cm}
    $\mathbf{E}$ should have the following effects:
    \begin{itemize}
      \item make the sample unusual enough to be an anomaly
      \item give it the shape and size of the anomaly kind we asked for
    \end{itemize}
    \vspace{0.1cm}
    Now, what is $\mathbf{E}$?
  \par}
\end{frame}

% ======================================================= 5. how we steer
\begin{frame}{How we steer}{Two forces, applied to the model's running guess
  $\hat x_0$ at every step}

  \begin{columns}[T]
    \begin{column}{0.615\textwidth}
      \begin{eqbox}{1. Be as unusual as normal data ever gets --- the \emph{shell}}
        \footnotesize
        \vspace{-0.75em}
        \begin{align*}
          h(x) &= -\tau \log \textstyle\sum_{r \in R}
                  e^{-\lVert e(x) - e(r)\rVert^2/\tau}
                  &&\text{\scriptsize how unusual $x$ is}\\[0em]
          Q &= \text{99th percentile of } h,\ \text{held-out nominal}
                  &&\text{}\\[0em]
          \text{shell} &= \bigl\{\, x \;:\; \lvert h(x) - Q\rvert \le \delta \,\bigr\}
                  &&\text{}
        \end{align*}
        \vspace{-1.6em}
      \end{eqbox}

      \vspace{0.02cm}
      \begin{eqbox}[good]{2. Have the shape and size we asked for}
        \footnotesize
        \vspace{-0.55em}
        \[
          \begin{gathered}
            \langle w,\, x\rangle \;=\; \textstyle\sum_t w[t]\,x[t]
            \;\longrightarrow\; a,
            \qquad \textstyle\sum_t w[t] = 0 \\[0.15em]
            \text{\scriptsize where $w$ encodes the shape and $a$ the size}
          \end{gathered}
        \]
        \vspace{-1.2em}
      \end{eqbox}

      \vspace{0.02cm}
      \begin{eqbox}[ink]{Together, the edit $\mathbf{E}$}
        \footnotesize
        \vspace{-0.9em}
        \[
          \mathbf{E} \;=\;
          \underbrace{\lambda\, \nabla_{\hat x_0}\bigl(h(\hat x_0) - Q\bigr)^2}
                     _{\text{\scriptsize onto the shell}}
          \;+\;
          \underbrace{\bigl(\langle w, \hat x_0\rangle - a\bigr)\, w
                      \,/\, \lVert w \rVert^2}
                     _{\text{\scriptsize to the asked-for shape}}
        \]
        \vspace{-1.12em}
      \end{eqbox}
    \end{column}

    \begin{column}{0.365\textwidth}
      \centering
      \includegraphics[width=\textwidth]{figs/shell.pdf}

      \vspace{-0.05cm}
      \raggedright
      {\footnotesize
       \textbf{In words.} Make the window as unusual as the most unusual 1\% of normal
       telemetry}
    \end{column}
  \end{columns}

  \vspace{0.2em}
  \noindent{\color{hair}\rule{\textwidth}{0.4pt}}\par
  \vspace{-0.4em}
  {\small\color{mute}
   $e(\cdot)$ is a trained encoder.\par}
\end{frame}

% ======================================================= 6. gallery
\begin{frame}{What the model produces}{three real windows and five generated ones, per kind}

  \begin{center}
    \includegraphics[width=0.97\textwidth]{figs/gallery.pdf}
  \end{center}

\end{frame}

% ======================================================= 7. realism
\begin{frame}{Realism and diversity}{Evaluation against other existing generators, on the
  two metrics the GenIAS paper reports}

  \begin{columns}[T]
    \begin{column}{0.53\textwidth}
      \begin{eqbox}{ARP --- realism}
        \footnotesize
        How close the \emph{real} anomalies are to their nearest generated window:
        \vspace{-0.4em}
        \[
          \mathrm{ARP} = \frac{1}{1 + \frac{1}{|A|}\sum_{a \in A}
            \min_{g \in G}\,\lVert \phi(a) - \phi(g)\rVert_2}
        \]
        \vspace{-1.5em}
      \end{eqbox}

      \vspace{0.06cm}
      \begin{eqbox}[good]{EDI --- diversity}
        \footnotesize
       Entropy-like term of the distribution of the generated windows:
        \vspace{-0.4em}
        \[
          \mathrm{EDI} = -\sum_{j=1}^{k} p_j \log p_j
        \]
        \vspace{-1.5em}
      \end{eqbox}

      \vspace{0.06cm}
      {}
    \end{column}

    \begin{column}{0.44\textwidth}
      \centering
      \tablestripe
      \small\renewcommand{\arraystretch}{1.5}
      \begin{tabular}{@{}l r r@{}}
        \toprule
          & ARP & EDI \\
        \midrule
        real anomalies \mt{(target)}               & 0.792 & 1.48 \\
        \textbf{ours}                              & \gd{0.509} & \gd{2.44} \\
        GenIAS                                     & 0.317 & 1.14 \\
        CutAddPaste                                & 0.323 & 1.27 \\
        handcrafted injection                      & 0.319 & 1.31 \\
        \bottomrule
      \end{tabular}

      \vspace{0.25cm}
      \raggedright
      {}
    \end{column}
  \end{columns}
\end{frame}

% ======================================================= 8. detection
\begin{frame}{Does it help a detector?}{$k$ real anomalies plus $k$ generated ones, against
  $k$ real alone. recall at 1\% false-alarm rate}

  \begin{columns}[T]
    \begin{column}{0.52\textwidth}
      \vspace{-0.1cm}
      \includegraphics[width=\textwidth]{figs/sweep.pdf}
    \end{column}

    \begin{column}{0.46\textwidth}
      \tablestripe
      \small\renewcommand{\arraystretch}{1.15}
      \begin{tabular}{@{}r r r r@{}}
        \toprule
        $k$ & real only & + generated & gain \\
        \midrule
          8 & 0.364 & 0.545 & \gd{+0.155\,*} \\
         16 & 0.455 & 0.591 & +0.164 \\
         32 & 0.364 & 0.636 & \gd{+0.200\,*} \\
         64 & 0.545 & 0.636 & \gd{+0.136\,*} \\
        128 & 0.591 & 0.727 & +0.118 \\
        256 & 0.591 & 0.682 & +0.109 \\
        \bottomrule
      \end{tabular}

      \vspace{0.1cm}
      {\scriptsize\color{mute}* 95\% confidence interval excludes zero}

      \vspace{0.2cm}
      {\small\color{good}\bfseries +0.200 event recall at $k=32$}\\
      {\footnotesize\color{good}4.4 more events of 22 detected}
    \end{column}
  \end{columns}

  \vspace{0.12cm}
  \begin{itemize}\setlength\itemsep{0.08em}\small
    \item \gd{Comparison against existing generators at $k=256$:\\ GenFSDiff 0.682 vs GenIAS 0.591, CutAddPaste 0.545.}
  \end{itemize}

  \vspace{0.04cm}
  {\scriptsize\color{mute}
   \textbf{The detector}: a 31k-parameter 1-D CNN on the raw window --- three convolutions
   with SiLU activations ($1\!\to\!32$ width 7, $32\!\to\!64$ width 5, $64\!\to\!64$ width
   5, each stride 2)
   \textbf{Held-out anomalies:} 22 anomaly events}
\end{frame}

% ======================================================= 9. summary
\begin{frame}{Summary, and what is not settled}

  \begin{columns}[T]
    \begin{column}{0.48\textwidth}
      \begin{eqbox}[good]{What we showed}
        \small
        \begin{itemize}\setlength\itemsep{0.4em}
          \item GenFSDiff generates synthetic anomalies of a
                \emph{chosen} kind and size, without retraining the denoiser.
          \item Adding them to training improves anomaly detection on
                ESA-ADB Mission~1.
        \end{itemize}
      \end{eqbox}
    \end{column}

    \begin{column}{0.48\textwidth}
      \begin{eqbox}[ink]{What is not settled}
        \small
        \begin{itemize}\setlength\itemsep{0.4em}
          \item \textbf{Continuous channels only.} Discrete telemetry
                (status flags, modes) would need e.g.\ discrete diffusion.
          \item \textbf{One channel at a time.} Anomalies that live in
                the relation between channels are not generated. Potential to include causality here?
        \end{itemize}
      \end{eqbox}
    \end{column}
  \end{columns}

\end{frame}
\end{document}
